\setcounter{subsection}{0}
\section{DISEÑO DE LA SOLUCIÓN}

El diseño de Kit Pagos Colombia se documentó en el SAD 2.0 (Equipo del Proyecto, 2026a). Este capítulo resume, sin reproducir ese documento, las decisiones que sostienen la propuesta del proyecto, y las relaciona con los requisitos y las restricciones que las originaron. La pregunta que guía el capítulo es por qué la arquitectura permite que un comercio integre varias pasarelas sin escribir el mapeo, la firma ni el flujo de consulta de cada una. La respuesta se organiza en seis partes: los criterios de diseño, la arquitectura hexagonal y el modelo de dominio, el diseño de los componentes, los patrones de integración, el tratamiento de la seguridad y los errores, y las tecnologías empleadas.

\subsection{Criterios y decisiones de diseño}

El diseño se derivó de los requisitos arquitectónicamente significativos (ASR, \textit{Architecturally Significant Requirements}) del SAD y de las restricciones del proyecto. Los ASR de prioridad alta son la extensibilidad ante nuevas pasarelas, la normalización uniforme de respuestas, la autenticidad de los webhooks, la protección de credenciales, la reproducción fiel de escenarios y la recuperación ante fallas transitorias. El manejo uniforme de errores y el bajo acoplamiento tienen prioridad media o baja. La Tabla~\ref{tab:asr-memoria} resume cada atributo, la decisión que lo resuelve y el escenario con el que se verifica. El SAD registra cada decisión como un ADR, con el problema que la motivó, la solución elegida y sus consecuencias.

\begin{table}[H]
  \centering
  \small
  \renewcommand{\arraystretch}{1.2}
  \begin{tabularx}{\textwidth}{|>{\raggedright\arraybackslash}p{2.9cm}|>{\raggedright\arraybackslash}X|>{\raggedright\arraybackslash}p{4.1cm}|}
    \hline
    \textbf{Atributo} & \textbf{Decisión de diseño} & \textbf{Escenario de verificación} \\ \hline
    Extensibilidad & Puerto \texttt{PaymentGatewayPort} y un adaptador por pasarela, seleccionado por fábrica. & Incorporar una pasarela nueva sin modificar el código del comercio ni el de las pasarelas existentes. \\ \hline
    Intercambiabilidad & Fachada \texttt{KitPagos} y configuración centralizada en \texttt{SdkConfigurator}. & Cambiar de pasarela modificando un solo campo de configuración. \\ \hline
    Normalización & \texttt{ResponseNormalizer} y tabla canónica de estados nativos. & Un mismo pago produce el mismo estado normalizado en las cuatro pasarelas. \\ \hline
    Seguridad & Tokenización en el navegador, verificación de firmas y enmascaramiento de secretos. & La tarjeta no llega al servidor del comercio; un webhook alterado se rechaza. \\ \hline
    Resiliencia & Reintentos con espera exponencial solo en operaciones idempotentes. & Una falla transitoria se recupera sin duplicar un cobro. \\ \hline
    Reproducibilidad & API de Simulación con escenarios deterministas y webhooks firmados. & Un caso de prueba produce el mismo desenlace en cada corrida. \\ \hline
  \end{tabularx}
  \caption{Atributos de calidad, decisiones de diseño y escenarios de verificación}
  \label{tab:asr-memoria}
\end{table}

Las restricciones acotaron el espacio de soluciones. El SDK se implementa exclusivamente en TypeScript sobre Node.js 20.19 o superior, y las pruebas con pasarelas reales dependen de sandboxes externos que el equipo no controla. El despliegue de la API de Simulación usa el plan gratuito de Render, y el proyecto dispone de cinco meses. Además, el framework no puede procesar ni almacenar datos de tarjeta, porque esa responsabilidad queda en los mecanismos certificados de cada proveedor. Estas restricciones explican decisiones que de otro modo parecerían arbitrarias, como el catálogo cerrado de URLs por ambiente o el almacenamiento efímero del simulador.

Tres principios guiaron la descomposición. El primero es la inversión de dependencias: el dominio no importa nada de los adaptadores, y son los adaptadores los que dependen de él. El segundo separa la configuración de la lógica, de modo que agregar una pasarela signifique para el comercio registrar sus credenciales en la configuración y no reescribir sus reglas de negocio. El tercero es que lo externo se pueda sustituir en las pruebas. Las llamadas HTTP se reemplazan sustituyendo \texttt{fetch}, y el reloj de la verificación de webhooks se puede inyectar, así que las suites corren sin red. Febrian et al. (2021) llegan a una idea cercana desde las líneas de producto de software. Tratan cada pasarela como una variación, encapsulan la comunicación con cada proveedor en un adaptador de API y reportan menos líneas de código que con el desarrollo convencional.

\subsection{Arquitectura hexagonal y modelo de dominio}

El núcleo del SDK sigue el estilo de arquitectura hexagonal, también llamado de puertos y adaptadores (Cockburn \& Garrido de Paz, 2024). El código vive en tres capas, \texttt{domain/}, \texttt{application/} e \texttt{infrastructure/}, y las dependencias apuntan siempre hacia el dominio. En el dominio están la entidad inmutable \texttt{Transaction}; los objetos de valor, entre ellos \texttt{Amount}, \texttt{Currency}, \texttt{PaymentResult}, \texttt{PendingRedirect}, \texttt{PseBank}, \texttt{RejectionReason} y \texttt{WebhookEvent}; el error \texttt{KitPagosError} y el servicio \texttt{WebhookVerifier}. No importa nada de las capas exteriores. La aplicación contiene el puerto \texttt{PaymentGatewayPort} y los servicios \texttt{ResponseNormalizer}, \texttt{RetryHandler} y \texttt{ErrorHandler}, que usan la fachada y los adaptadores. La infraestructura implementa los detalles: los cuatro adaptadores, que llaman a \texttt{fetch} con un tiempo límite, la configuración, la fábrica, la fachada y el enmascaramiento de credenciales.

La regla de dependencia se concreta en el puerto de salida \texttt{PaymentGatewayPort}, que declara las operaciones \texttt{createPayment}, \texttt{getStatus}, \texttt{getPseBanks} y \texttt{verifySignature}. Cada pasarela implementa ese puerto en su propio adaptador. La fachada \texttt{KitPagos} depende del puerto y de los tipos del dominio, nunca de una clase concreta de proveedor; la selección del adaptador ocurre en \texttt{GatewayFactory}. Esta estructura es la que permite que el código del comercio ignore el vocabulario de cada pasarela y que el cambio de proveedor sea una operación de configuración (Martin, 2017). Dentro del SDK, en cambio, una pasarela nueva sí extiende el dominio: suma un valor al enumerado \texttt{Gateway}, su tabla de estados nativos y su manejador de webhooks, además de su adaptador y su normalizador. Lo que no cambia es el código de las pasarelas existentes. El modelo de dominio sigue el principio de lenguaje ubicuo del diseño guiado por el dominio (DDD, \textit{Domain-Driven Design}) (Evans, 2003). Cada término del negocio tiene un único identificador en el código, como \texttt{Transaction}, \texttt{Amount}, \texttt{OrderReference}, \texttt{PseBank} y \texttt{RejectionReason}. El glosario del dominio fija la correspondencia de cada uno con el vocabulario en español (Artefacto~\ref{anx:arquitectura}).

Las Figuras~\ref{fig:memoria-contenedores} y~\ref{fig:memoria-componentes} presentan las vistas C4 de contenedores y de componentes del SDK. Son los mismos diagramas del SAD, elaborados con el modelo C4 (Brown, s.~f.), que describe un sistema en cuatro niveles de detalle: contexto, contenedores, componentes y código. Reutilizarlos hace que la memoria y el documento de arquitectura describan el mismo sistema sin mantener dos versiones de cada figura.

\begin{figure}[H]
  \centering
  \includegraphics[width=\textwidth,height=0.78\textheight,keepaspectratio]{"../../architecture/SDK/Container Diagram - C4.pdf"}
  \caption{Diagrama de contenedores de Kit Pagos Colombia}
  \label{fig:memoria-contenedores}
\end{figure}

\begin{figure}[H]
  \centering
  \includegraphics[width=\textwidth,height=0.78\textheight,keepaspectratio]{"../../architecture/SDK/Component Diagram - C4.pdf"}
  \caption{Diagrama de componentes del SDK}
  \label{fig:memoria-componentes}
\end{figure}

\subsection{Diseño de los componentes}

El SDK es el contenedor de mayor complejidad. Su parte de servidor agrupa los componentes que describe la Tabla~\ref{tab:componentes-memoria}. El módulo de navegador se distribuye en el mismo paquete, en la subruta \texttt{kit-pagos-colombia/browser}, y no depende de módulos de Node.js. La fachada \texttt{KitPagos} expone las operaciones del comercio y delega en la configuración y en la fábrica. \texttt{SdkConfigurator} recibe la pasarela, las credenciales, el ambiente, la URL base, los reintentos, la tolerancia de webhooks y el tiempo límite. Exige la pasarela y las credenciales, valida el ambiente y el tiempo límite, y \texttt{GatewayFactory} instancia el adaptador correspondiente. \texttt{ResponseNormalizer} traduce los estados nativos a los seis estados canónicos (\texttt{APPROVED}, \texttt{DECLINED}, \texttt{PENDING}, \texttt{EXPIRED}, \texttt{VOIDED} y \texttt{ERROR}) y conserva el valor original en \texttt{rawStatus}. \texttt{WebhookVerifier} comprueba la firma de cada proveedor; \texttt{RetryHandler} aplica la política de reintentos; \texttt{ErrorHandler} convierte las fallas en \texttt{KitPagosError}.

\begin{table}[H]
  \centering
  \small
  \renewcommand{\arraystretch}{1.2}
  \begin{tabularx}{\textwidth}{|>{\raggedright\arraybackslash}p{3.6cm}|>{\raggedright\arraybackslash}X|}
    \hline
    \textbf{Componente} & \textbf{Responsabilidad} \\ \hline
    \texttt{KitPagos} & Fachada que expone creación, consulta, bancos PSE y verificación de webhooks al comercio. \\ \hline
    \texttt{SdkConfigurator} y \texttt{GatewayFactory} & Reciben y validan las opciones y resuelven el adaptador de la pasarela configurada. \\ \hline
    \texttt{PaymentGatewayPort} y adaptadores & Contrato común y traducción del contrato del SDK al de Wompi, Mercado Pago, Kushki y Rapyd. \\ \hline
    \texttt{ResponseNormalizer} & Traduce estados, motivos de rechazo y montos a los tipos del dominio. \\ \hline
    \texttt{WebhookVerifier} y \texttt{ErrorHandler} & Verifican firmas y convierten fallas en errores tipados, sin credenciales a la vista. \\ \hline
    \texttt{RetryHandler} & Reintenta solo operaciones idempotentes, con una espera que crece de forma exponencial y una variación aleatoria. \\ \hline
    Módulo de navegador & Tokeniza la tarjeta en las pasarelas que lo permiten, sin exponerla al servidor del comercio. \\ \hline
  \end{tabularx}
  \caption{Componentes principales del SDK y su responsabilidad}
  \label{tab:componentes-memoria}
\end{table}

La API de Simulación es un servicio Fastify que atiende dos grupos de rutas HTTP en un solo proceso y un solo despliegue. La primera, \texttt{/v1/sim}, reproduce el comportamiento nativo de las cuatro pasarelas para que las pruebas no dependan de sus sandboxes. Recibe la petición con el formato del proveedor, resuelve un escenario y responde con el mismo vocabulario que él. La segunda, \texttt{/v1/api}, ofrece por HTTP las operaciones de la fachada \texttt{KitPagos} y está pensada para clientes que no instalan el SDK de servidor, como una aplicación de navegador o una herramienta de pruebas. Los dos grupos comparten proceso, puerto y credenciales, pero cumplen funciones distintas: \texttt{/v1/sim} simula a la pasarela y guarda sus registros, y \texttt{/v1/api} consume el SDK, que en modo simulación llama por HTTP a \texttt{/v1/sim}. Esta distinción importa al leer los resultados: una prueba sobre \texttt{/v1/api} pone a prueba el SDK completo, mientras que una prueba sobre \texttt{/v1/sim} solo prueba el simulador.

Para elegir el escenario, el simulador mira en orden la cabecera \texttt{x-simulator-scenario}, el dato de prueba nativo de la pasarela y un monto reservado; si ninguno aplica, responde \texttt{APPROVED}. Los almacenes en memoria guardan los registros por identificador nativo y las máquinas de estado modelan las transiciones válidas de cada proveedor. Los webhooks se generan con la firma de la pasarela correspondiente y se envían a una URL configurada en el servidor, nunca a una dirección tomada del cuerpo de la petición. La documentación OpenAPI 3.0.3 se publica en \texttt{/docs}, \texttt{/docs/json} y \texttt{/docs/yaml}.

\subsection{Patrones de integración}

El sistema implementa dos patrones de integración completos y un tercero en parte, como resume la Tabla~\ref{tab:patrones-memoria}. En P1, el servidor del comercio usa el SDK y este se comunica con la pasarela. En P2, una aplicación de navegador consume la API REST \texttt{/v1/api}, que a su vez usa el SDK; en simulación, el SDK llama por HTTP a \texttt{/v1/sim} dentro del mismo servicio. P3 es un SDK unificado en el navegador. De él existe el módulo \texttt{KitPagosBrowser}, que tokeniza la tarjeta en Wompi, Mercado Pago y Kushki y que el caso de prueba (CP) CP-05 del plan de pruebas verifica. El cobro sigue pasando por un servidor, el del comercio en P1 o \texttt{/v1/api} en P2, y Rapyd no tokeniza en el navegador: cobra la tarjeta en su página alojada.

\begin{table}[H]
  \centering
  \small
  \renewcommand{\arraystretch}{1.2}
  \begin{tabularx}{\textwidth}{|>{\centering\arraybackslash}p{1.25cm}|>{\raggedright\arraybackslash}X|>{\centering\arraybackslash}p{2.2cm}|}
    \hline
    \textbf{Patrón} & \textbf{Flujo} & \textbf{Estado} \\ \hline
    P1 & Servidor del comercio $\rightarrow$ SDK $\rightarrow$ pasarela & Implementado \\ \hline
    P2 & Navegador $\rightarrow$ \texttt{/v1/api} $\rightarrow$ SDK $\rightarrow$ pasarela (en simulación, \texttt{/v1/sim}) & Implementado \\ \hline
    P3 & Navegador $\rightarrow$ SDK unificado en el navegador $\rightarrow$ pasarela (solo la tokenización, sin Rapyd) & En parte \\ \hline
  \end{tabularx}
  \caption{Patrones de integración soportados y explorados}
  \label{tab:patrones-memoria}
\end{table}

Los patrones se apoyan en ideas clásicas de integración de sistemas: el adaptador traduce el contrato, el enrutador de mensajes dirige la petición al proveedor correcto y el normalizador convierte la respuesta a un formato común (Hohpe \& Woolf, 2003). Hisyam y Manuaba (2022) estudiaron un caso parecido: integraron dos pasarelas indonesias, Midtrans y Xendit, en el servidor de una plataforma de comercio con varios vendedores (\textit{marketplace}) para dividir un mismo pago entre billeteras electrónicas de proveedores distintos. Kit Pagos Colombia atiende otro caso, el del comercio que cobra con una pasarela a la vez y quiere poder cambiarla, y no aborda el pago dividido.

\subsection{Seguridad y tratamiento de errores}

La seguridad se apoya en delegar lo sensible. El framework no procesa ni almacena datos de tarjeta. En Wompi, Mercado Pago y Kushki, el número viaja del navegador a la ruta de tokenización de la pasarela, o del simulador en las pruebas, y el servidor del comercio recibe solo un token opaco. En Rapyd, la tarjeta se captura en la página alojada del proveedor. Esta separación saca los datos de la tarjeta del servidor del comercio, aunque no lo deja fuera de PCI DSS. Si la página de pago es del comercio y tokeniza con JavaScript, el comercio queda en el cuestionario SAQ A-EP (PCI Security Standards Council, 2024). Cada ambiente tiene su catálogo de URLs. En el módulo de navegador el catálogo es cerrado; el SDK de servidor solo lo deja de lado si se configura \texttt{baseUrl} de forma explícita. Los errores ocultan las llaves antes de propagarse. La verificación de webhooks usa el algoritmo de firma de cada proveedor y una tolerancia temporal configurable; sin esa verificación, un tercero podría alterar el estado de una orden con una notificación falsa. La seguridad de las pasarelas de pago también se ha revisado en la literatura. Olanrewaju et al. (2017) comparan las pasarelas más usadas en seguridad, costo y soporte, y recogen las limitaciones de seguridad que conviene mitigar.

El tratamiento de errores distingue lo transitorio de lo definitivo. \texttt{RetryHandler} reintenta solo lecturas idempotentes, es decir, consultas de estado y listas de bancos, con una espera que crece de forma exponencial, una variación aleatoria y un número máximo de intentos. La creación de un pago no se reintenta, porque un reintento ciego podría cobrar dos veces. Cada adaptador traduce sus fallas a \texttt{KitPagosError} con \texttt{ErrorHandler}, con su código, la respuesta original y un mensaje sin credenciales. Si una falla transitoria persiste después de los reintentos, \texttt{RetryHandler} la envuelve en \texttt{MAX\_RETRIES\_EXCEEDED} con el último error como causa. Este patrón de aislar fallas y evitar reintentos no idempotentes es el que la literatura de resiliencia recomienda para sistemas que dependen de servicios externos (Nygard, 2018). También responde a un riesgo documentado en la integración de APIs de pago. En el estudio de Aué et al. (2018), la mayoría de las fallas se debía a datos inválidos o incompletos en las peticiones de quien integraba.

Los montos se representan como una cadena decimal exacta, y no como un número de punto flotante, para no perder centavos en la conversión. El sistema respeta los códigos de moneda de ISO~4217 y evita el redondeo de la aritmética binaria de punto flotante de IEEE~754 (International Organization for Standardization, 2015; Institute of Electrical and Electronics Engineers, 2019). En el simulador, el monto también funciona como regla de escenario cuando cae en un valor reservado; esa convención es una herramienta de prueba y no una regla del framework. El capítulo VII declara qué asegura el simulador y qué solo puede verificarse con una pasarela real.

\subsection{Tecnologías y herramientas}

La Tabla~\ref{tab:tecnologias-memoria} resume las tecnologías y su justificación. La elección de Node.js y TypeScript busca que el SDK y el simulador compartan lenguaje, tipos y herramientas de prueba, y que el paquete pueda publicarse en npm sin un entorno de ejecución adicional (Node.js, s.~f.). Vitest se usa en el prototipo C, y Jest, en las suites del SDK y del simulador. Fastify sirve la API y OpenAPI documenta sus rutas; Postman y Newman ejecutan la colección de extremo a extremo, Render aloja el despliegue y GitHub Actions ejecuta la integración continua. Las métricas CK se calculan con un script propio que usa la biblioteca \texttt{ts-morph}, y la carga se mide con Autocannon. El paquete se publica en npm como \texttt{kit-pagos-colombia}; la versión evaluada es la 0.7.0-rc.2, publicada con la etiqueta \texttt{next} (Kit Pagos Colombia, 2026).

\begin{table}[H]
  \centering
  \small
  \renewcommand{\arraystretch}{1.2}
  \begin{tabularx}{\textwidth}{|>{\raggedright\arraybackslash}p{2.9cm}|>{\raggedright\arraybackslash}p{4.7cm}|>{\raggedright\arraybackslash}X|}
    \hline
    \textbf{Tecnología} & \textbf{Uso} & \textbf{Justificación} \\ \hline
    Node.js 20 y TypeScript & SDK, simulador y prototipos; el SDK exige Node.js 20.19 o superior & Un solo lenguaje y tipos compartidos entre servidor y pruebas. \\ \hline
    Jest y Vitest & Pruebas unitarias, de integración y de contrato & Cobertura, sustitución de red y ejecución reproducible. \\ \hline
    Fastify y OpenAPI 3.0.3 & API de Simulación y documentación & Rendimiento, esquemas y documentación autoexplorable. \\ \hline
    Postman y Newman & Colección de extremo a extremo & Pruebas ejecutables en local y en Render sin instalar el SDK. \\ \hline
    Render y GitHub Actions & Despliegue e integración continua & URL pública para la evaluación y verificación automática en cada cambio. \\ \hline
    ts-morph y Autocannon & Métricas CK y carga & Medición reproducible de complejidad y comportamiento bajo concurrencia. \\ \hline
  \end{tabularx}
  \caption{Tecnologías y herramientas del framework}
  \label{tab:tecnologias-memoria}
\end{table}

La comparación del capítulo VII depende de que el diseño separe la traducción del contrato de cada pasarela, la normalización del resultado y la orquestación del caso de uso. Con esa separación se puede medir cuánto código escribe el comercio cuando el framework hace la traducción y cuánto cuando integra cada pasarela por su cuenta.
